% ===== PRÉAMBULE CANONIQUE — à coller VERBATIM en tête de CHAQUE .tex =====
\documentclass[11pt,a4paper]{article}
\usepackage[utf8]{inputenc}\usepackage[T1]{fontenc}\usepackage{lmodern}
\usepackage[french]{babel}
\usepackage{amsmath,amssymb,mathtools}\usepackage{icomma}
\usepackage{siunitx}\sisetup{locale=FR}  % \qty{}{}, \si{}, \ang{}
\usepackage{xcolor}\usepackage{colortbl}
\usepackage{tikz}\usepackage{pgfplots}\pgfplotsset{compat=1.18}
\usepackage[siunitx,europeanresistors]{circuitikz} % circuits (élec.)
\usepackage[most]{tcolorbox}\usepackage{enumitem}\usepackage{array,booktabs}
\usepackage[a4paper,margin=2.2cm]{geometry}
\usetikzlibrary{arrows.meta,calc,angles,quotes,positioning,patterns,%
  backgrounds,shapes.geometric,decorations.pathmorphing,decorations.markings,intersections}
\definecolor{primary}{RGB}{222,49,99}\definecolor{primarydark}{RGB}{150,22,55}
\definecolor{defcol}{RGB}{0,114,178}\definecolor{methcol}{RGB}{52,92,148}
\definecolor{excol}{RGB}{0,135,90}\definecolor{attncol}{RGB}{200,40,55}
\tcbset{boxbase/.style={enhanced,breakable,boxrule=0.9pt,arc=2.5pt,
  left=3mm,right=3mm,top=2.3mm,bottom=2.3mm,fonttitle=\bfseries,coltitle=white}}
% --- boîtes TUTORIEL (noms EXACTS, ne pas renommer) ---
\newtcolorbox{cle}[1][]{boxbase,colback=defcol!6,colframe=defcol,
  title={\textsc{À retenir}\ifx\relax#1\relax\relax\else\textmd{ : \itshape #1}\fi}}
\newtcolorbox{methode}[1][]{boxbase,colback=methcol!7,colframe=methcol,sharp corners,
  title={\textsc{Méthode}\ifx\relax#1\relax\relax\else\textmd{ --- \itshape #1}\fi}}
\newtcolorbox{exemplebox}[1][]{boxbase,colback=excol!5,colframe=excol,
  title={\textsc{Exemple}\ifx\relax#1\relax\relax\else\textmd{ : \itshape #1}\fi}}
\newtcolorbox{piege}[1][]{boxbase,colback=attncol!6,colframe=attncol,
  title={\textsc{Piège}\ifx\relax#1\relax\relax\else\textmd{ : \itshape #1}\fi}}
\newtcolorbox{remarque}[1][]{boxbase,colback=black!4,colframe=black!45,
  title={\textsc{Remarque}\ifx\relax#1\relax\relax\else\textmd{ : \itshape #1}\fi}}
% --- boîtes EXERCICE / CORRIGÉ (noms EXACTS, auto-comptées) ---
\newtcolorbox[auto counter]{exo}[1][]{boxbase,colback=primary!4,colframe=primary,
  title={\textsc{Exercice}~\thetcbcounter\ifx\relax#1\relax\relax\else\textmd{ --- \itshape #1}\fi}}
\newtcolorbox[auto counter]{corrige}[1][]{boxbase,colback=excol!4,colframe=excol,
  title={\textsc{Corrigé}~\thetcbcounter\ifx\relax#1\relax\relax\else\textmd{ --- \itshape #1}\fi}}
\newcommand{\vv}[1]{\vec{#1}}\newcommand{\ds}{\displaystyle}
\newcommand{\R}{\mathbb{R}}\newcommand{\N}{\mathbb{N}}\newcommand{\Z}{\mathbb{Z}}
% (un \usepackage{fancyhdr}+en-tête est possible ; il est ignoré côté web)
% =========================================================================
\begin{document}
\begin{corrige}[la boussole déviée par un fil]
Les deux champs horizontaux sont perpendiculaires : on compose par Pythagore. Avec
$B_f=\qty{3.46e-5}{\tesla}=2\sqrt{3}\times10^{-5}$ :
\begin{enumerate}[label=\alph*)]
  \item $B_{\text{tot}}=\sqrt{B_H^2+B_f^2}=\sqrt{(2)^2+(2\sqrt3)^2}\times10^{-5}
        =\sqrt{4+12}\times10^{-5}=\sqrt{16}\times10^{-5}=\qty{4e-5}{\tesla}.$
  \item L'aiguille s'aligne sur $\vv{B}_{\text{tot}}$ :
        \[ \tan\alpha=\frac{B_f}{B_H}=\frac{3{,}46}{2}=\sqrt{3}\;\Rightarrow\;\alpha=\ang{60}. \]
        L'aiguille tourne de \ang{60} \textbf{vers l'Est}. C'est précisément l'expérience d'Ørsted :
        le courant « détourne » la boussole d'autant plus que $B_f$ est grand devant $B_H$.
\end{enumerate}
\end{corrige}

\begin{corrige}[deux fils de courants opposés : champ au milieu]
\begin{enumerate}[label=\alph*)]
  \item \textbf{Fil 1} ($\odot$, à gauche de $M$) : main droite $\Rightarrow$ $\vv{B}_1$ dirigé
        \textbf{vers le haut} en $M$. \textbf{Fil 2} ($\otimes$, à droite de $M$) : courant entrant,
        les lignes tournent en sens horaire ; à gauche du fil 2 (donc en $M$) le champ $\vv{B}_2$ est
        aussi dirigé \textbf{vers le haut}. Les deux champs sont \textbf{de même sens}.
  \item Ils s'\textbf{ajoutent} :
        \[ B_{\text{tot}}=B_1+B_2=4\cdot10^{-5}+4\cdot10^{-5}=\qty{8e-5}{\tesla}\ \text{(vers le haut).} \]
        \textbf{À retenir} : au milieu de deux fils, courants de \emph{même} sens $\Rightarrow$ champ
        \textbf{nul} ; courants \emph{opposés} $\Rightarrow$ champ \textbf{maximal} (les contributions
        s'ajoutent). C'est l'inverse de l'intuition !
\end{enumerate}
\end{corrige}

\begin{corrige}[superposition à \ang{60} : passer par les composantes]
\begin{enumerate}[label=\alph*)]
  \item Sur l'axe horizontal ($x$) et vertical ($y$), avec $B=\qty{4e-4}{\tesla}$ :
        \[ \vv{B}_1=(B,\;0),\qquad \vv{B}_2=(B\cos\ang{60},\;B\sin\ang{60})=(2\cdot10^{-4},\;3{,}46\cdot10^{-4}). \]
        Composantes du résultant :
        \[ R_x=B+B\cos\ang{60}=4\cdot10^{-4}+2\cdot10^{-4}=\qty{6e-4}{\tesla},\quad
           R_y=B\sin\ang{60}=\qty{3.46e-4}{\tesla}. \]
  \item Norme :
        \[ B_{\text{tot}}=\sqrt{R_x^2+R_y^2}=\sqrt{6^2+3{,}46^2}\times10^{-4}
           \approx\sqrt{48}\times10^{-4}\approx\qty{6.9e-4}{\tesla}. \]
        Angle avec l'horizontale : $\tan\varphi=\dfrac{R_y}{R_x}=\dfrac{3{,}46}{6}=0{,}577
        \Rightarrow \varphi=\ang{30}$. La résultante \textbf{bissecte} l'angle des deux champs
        (normal, puisqu'ils ont la même norme).
\end{enumerate}
\end{corrige}

\begin{corrige}[annuler un champ avec un fil]
\begin{enumerate}[label=\alph*)]
  \item Pour que $\vv{B}_A+\vv{B}_f=\vv{0}$, il faut $\vv{B}_f=-\vv{B}_A$ : même intensité
        $\qty{6e-5}{\tesla}$, même direction (horizontale), mais sens \textbf{opposé}, donc dirigé
        \textbf{vers la gauche}.
  \item Le point $M$ est \emph{au-dessous} du fil. On veut qu'en ce point le champ pointe vers la
        gauche. Par la règle de la main droite, cela impose un courant \textbf{entrant} dans la feuille
        ($\otimes$) : pouce enfoncé dans la feuille, les doigts tournent en sens horaire, et sous le
        fil la tangente est bien dirigée vers la gauche. Avec ce sens et une intensité réglée pour
        $B_f=\qty{6e-5}{\tesla}$, le champ total est nul en $M$.
\end{enumerate}
\end{corrige}

\end{document}
